\section{Introduction}
\label{sec:introduction}

Network intrusion detection systems (NIDS) identify malicious activity in network traffic.
For attacks spanning multiple flows, evidence is distributed across requests, responses, and communications among hosts, so the security significance of an individual interaction depends on the surrounding behavior.
A failed request can reflect a configuration error, while repeated attempts by the same client across multiple resources, successive server rejections, and subsequent changes in responses can provide additional evidence of probing or intrusion.
Detecting such behavior requires NIDS to combine the operational meaning of related communications while processing a continuous stream of traffic.

Machine learning methods extract traffic representations from packet bytes, protocol structures, and relationships among flows.
NetMamba learns communication patterns from packet bytes through pretraining~\cite{R06}, while the structure-aware method of Yu et al. preserves field boundaries and protocol hierarchies~\cite{E01}.
Cross-flow methods extend the evidence available for detection by incorporating temporal relationships, shared endpoints, and statistical similarity.
Tracegram aggregates flows with temporal context~\cite{R01}, and the framework of Huang et al. uses relation-aware attention to incorporate cross-flow structure into target-flow classification~\cite{R07}.

The operational meaning of communications remains largely implicit in these representations.
Explicit alignment of operations, accessed resources, and response outcomes provides a common basis for interpreting messages from different protocols and combining them into behavioral evidence.
Semantic methods make operational semantics explicit: the BACnet threat detector uses large language models (LLMs) to interpret protocol operations with domain knowledge~\cite{R12}, while MISSION correlates MQTT communications and constructs process models~\cite{R13}.
Their mechanisms for interpreting operations and correlating interactions are defined for specific protocols.
For heterogeneous traffic, the first challenge is to align protocol-specific operations and organize related interactions into bounded evidence that preserves operation order and response outcomes for target-flow classification.

Cross-flow semantic analysis also introduces processing costs that depend on context size, decision procedure, and invocation count.
Generative LLM detectors encode the input and generate an output sequence to obtain a decision; From Flows to Words follows this procedure with grammar-constrained structured output~\cite{R14}.
Submitting cross-flow context independently for each target flow can additionally repeat the processing of shared evidence.
Jev-IDS applies typed decisions to one serialized flow per request and evaluates detection quality, latency, and cost~\cite{R02}.
Applying semantic decisions to the full traffic stream requires candidate selection for two tasks: verifying alerts produced by a fast detector and identifying attacks supported jointly by multiple flows.
These tasks target precision and recall, respectively, and require different selection criteria.
The second challenge is to allocate semantic inference to both tasks while controlling context size and invocation costs.

This paper presents a network intrusion detection framework that combines a two-layer semantic representation with two parallel decision paths based on Jev.
The first layer normalizes heterogeneous protocol messages into operations, resources, response outcomes, and endpoint roles to establish message-level operational semantics.
The second layer constructs cross-flow behavioral semantics around a target flow, using endpoint and temporal associations to organize communication history, client activity, server responses, and related events in temporal order.
Deterministic computation extracts counts, rates, and response differences, while template folding compresses repeated messages and retains meaningful changes.
The resulting context presents behavioral facts relevant to the target flow within a fixed input budget and maintains links to the source packets and flows.

The two decision paths address alert verification and cross-flow attack discovery in parallel.
In the alert-verification path, a fast detector scores every flow, and Jev reassesses the resulting alert set to reduce false positives and improve precision.
Independently, the cross-flow discovery path applies selective routing across the full traffic stream, using association cues to select target flows whose related communications may jointly indicate an attack.
The two semantic layers organize this cross-flow evidence into a compact context for Jev to detect attacks that are difficult to identify from individual flows alone, thereby improving recall.
Jev evaluates typed questions against the supplied state and directly returns decision probabilities~\cite{R11}.
The framework controls semantic processing costs through bounded evidence construction, direct probability outputs, and candidate selection specific to each path.

Experiments on xxx show that the complete framework improves precision and recall by xxx and xxx percentage points, respectively, over the fast detector alone.
Ablation experiments show that alert verification reduces false positives by xxx\%, while cross-flow discovery increases the detection rate for attacks requiring related-flow evidence by xxx percentage points.
Compared with the xxx generative LLM under matched input evidence, detection operating points, and concurrency, the framework reduces end-to-end processing latency by xxx\%.

The main contributions are as follows:
\begin{enumerate}
    \item \textbf{A two-layer semantic representation for target-flow classification.}
    Message-level operational semantics align operations, resources, response outcomes, and endpoint roles across protocols.
    Cross-flow behavioral semantics organize related interactions and their temporal order around a target flow, with deterministic statistics and template folding producing bounded evidence traceable to the original traffic.
    \item \textbf{A dual-path semantic decision framework for precision and recall.}
    Alert verification reassesses the fast detector's alerts, while cross-flow discovery independently selects candidates whose related communications may jointly indicate an attack.
    The two paths combine Jev's typed probability outputs with task-specific candidate selection and compact semantic contexts to support false-positive correction and additional attack detection within a controlled processing budget.
\end{enumerate}
